% =====================================================================
% Resultados del Experimento 21 — Agentes especializados y Agent Skills
% Datos: Experimentos/experiments/experimento_21_agentskills/results/
% Corpus de evaluación: 1.313 noticias anotadas en las cinco variables
% =====================================================================

Esta sección presenta los resultados de la evaluación de la arquitectura de
agentes especializados descrita en el Capítulo~\ref{implementation}. Todas las
cifras se calculan sobre el subconjunto de \textbf{1.313 noticias} del corpus
IMIO que cuentan con anotación experta simultánea en las cinco variables de
lenguaje sexista consideradas (V25, V26, V30, V33 y V35).

\subsection{Diseño experimental y condiciones controladas}
\label{subsec:iris_diseno}

La aportación central de este trabajo es la arquitectura de agentes
especializados con \textit{Agent Skills}, descrita en el
Capítulo~\ref{implementation} y denotada aquí como \textbf{nivel B1}. Para
evaluarla con honestidad no basta con medir su rendimiento en aislado: hace falta
una \textbf{condición de control} que reciba exactamente la misma metodología por
una vía más simple. Esa condición es el \textbf{nivel B0}. Ambos niveles comparten
corpus, modelos, número de llamadas (cinco por artículo, una por variable) y
formato de salida. Solo cambia \emph{cómo llega la metodología al modelo}:

\begin{itemize}
    \item \textbf{Nivel B1 (\textit{Agent Skills}, sistema propuesto).} El agente
    recibe únicamente los metadatos (\texttt{name} y \texttt{description}) de las
    \textit{skills} disponibles y carga bajo demanda el cuerpo de la que necesita
    (\textit{progressive disclosure}), pudiendo además consultar las guías
    expertas y verificar sus evidencias mediante herramientas.
    \item \textbf{Nivel B0 (\textit{baseline}, control).} La metodología de cada
    variable se inyecta íntegramente en el \textit{prompt}. El modelo emite su
    veredicto en una única llamada, sin herramientas.
\end{itemize}

El papel del control es aislar la aportación de la arquitectura. Que B0 y B1
compartan todo salvo el mecanismo de entrega evita el sesgo de comparación
presente en trabajos previos del proyecto, donde el nivel con \textit{skills}
empleaba una única llamada conjunta frente a cinco llamadas independientes del
\textit{baseline}, de modo que cualquier diferencia observada resultaba
inatribuible. Para garantizar la
comparabilidad entre proveedores se desactivó el razonamiento explícito en
todos los modelos (\texttt{reasoning\_effort} en OpenAI,
\texttt{thinking\_budget} en Google) y se empleó el mismo fichero de
definiciones de variables en régimen estricto.

La Tabla~\ref{tab:iris-prevalencia} recoge la prevalencia de la clase positiva
en la anotación experta, dato imprescindible para interpretar correctamente las
métricas posteriores.

\begin{table}[H]
    \centering
    \begingroup
    \small
    \ttabbox[\FBwidth]{
        \caption{Prevalencia de la clase positiva en la anotación experta ($N=1.313$)}
        \label{tab:iris-prevalencia}
    }{
        \begin{tabularx}{\textwidth}{
            X
            >{\centering\arraybackslash}X
            >{\centering\arraybackslash}X
        }
            \hline
            \textbf{Variable} & \textbf{Código} & \textbf{Prevalencia «Sí»} \\
            \hline
            \texttt{lenguaje\_sexista}         & V25 & 0,810 \\
            \texttt{masc\_generico}            & V26 & 0,810 \\
            \texttt{sexismo\_discurso}         & V30 & 0,428 \\
            \texttt{asimetria\_mujer\_hombre}  & V33 & 0,062 \\
            \texttt{denominacion\_sexualizada} & V35 & 0,100 \\
            \hline
        \end{tabularx}
    }
    \endgroup
\end{table}

Dos rasgos condicionan toda la interpretación posterior. En primer lugar, V25 y
V26 presentan una prevalencia muy elevada (81\,\%), lo que implica que un
clasificador trivial que respondiese siempre «Sí» alcanzaría una exactitud de
0,810. En segundo lugar, V33 y V35 están fuertemente desbalanceadas (6,2\,\% y
10,0\,\% de positivos), régimen en el que el estadístico $\kappa$ resulta
especialmente severo.

\subsection{Resultados del núcleo: \textit{baseline} frente a Agent Skills}
\label{subsec:iris_core}

La Tabla~\ref{tab:iris-b0b1} presenta la comparación agregada entre ambos
niveles para los tres proveedores evaluados.

\begin{table}[H]
    \centering
    \begingroup
    \small
    \ttabbox[\FBwidth]{
        \caption{Comparación B0 (\textit{baseline}) frente a B1 (Agent Skills). Medias sobre las cinco variables, $N=1.313$}
        \label{tab:iris-b0b1}
    }{
        \begin{tabularx}{\textwidth}{
            X
            >{\centering\arraybackslash}X
            >{\centering\arraybackslash}X
            >{\columncolor{gray!15}\centering\arraybackslash}X
            >{\centering\arraybackslash}X
            >{\centering\arraybackslash}X
        }
            \hline
            \textbf{Modelo} & \textbf{Nivel} & \textbf{Exactitud} & \textbf{$\kappa$} & \textbf{F1 macro} & \textbf{Coste (USD)} \\
            \hline
            \multirow{3}{*}{\texttt{gemini-3.1-flash-lite}}
              & B0 & 0,600 & $+0,026$ & 0,388 & 5,64 \\
              & B1 & 0,661 & $\mathbf{+0,066}$ & 0,443 & 13,16 \\
              & $\Delta$ & $+0,061$ & $+0,040$ & $+0,055$ & $2,3\times$ \\
            \hline
            \multirow{3}{*}{\texttt{gpt-4o-mini}}
              & B0 & 0,564 & $+0,033$ & 0,410 & 2,52 \\
              & B1 & 0,612 & $+0,014$ & 0,438 & 6,89 \\
              & $\Delta$ & $+0,048$ & $-0,019$ & $+0,029$ & $2,7\times$ \\
            \hline
            \multirow{3}{*}{\texttt{gpt-5.4-nano}}
              & B0 & 0,571 & $+0,052$ & 0,430 & 5,67 \\
              & B1 & 0,569 & $+0,004$ & 0,364 & 10,86 \\
              & $\Delta$ & $-0,003$ & $-0,048$ & $-0,066$ & $1,9\times$ \\
            \hline
            \multirow{2}{*}{\textbf{Media}}
              & \textbf{B0} & 0,578 & $\mathbf{+0,037}$ & 0,409 & 13,83 \\
              & \textbf{B1} & \textbf{0,614} & $+0,028$ & 0,415 & 30,91 \\
            \hline
        \end{tabularx}
    }
    \endgroup
\end{table}

El hallazgo más matizado, y el más informativo, es que \textbf{el efecto de las
\textit{Agent Skills} depende del modelo}. En \texttt{gemini-3.1-flash-lite}, la
arquitectura propuesta mejora el acuerdo de forma clara: el coeficiente $\kappa$
pasa de $+0,026$ a $+0,066$ ($\Delta\kappa = +0,040$), y lo hace acompañado de
subidas simultáneas en exactitud y F1, señal de una ganancia genuina y no de un
artefacto de la distribución (Sección~\ref{subsec:iris_divergencia}). En los dos
modelos de OpenAI, en cambio, las \textit{skills} no ayudan e incluso perjudican
($\Delta\kappa = -0,019$ en \texttt{gpt-4o-mini} y $-0,048$ en
\texttt{gpt-5.4-nano}).

Esta dependencia del modelo tiene una explicación en el comportamiento del agente
(Sección~\ref{subsec:iris_lift}): las \textit{Agent Skills} aportan cuando el
modelo aprovecha de verdad la carga bajo demanda y las herramientas, y estorban
cuando el modelo sobrecarga el catálogo o infrautiliza las acciones. La
arquitectura, por tanto, no es útil ni inútil en abstracto: su valor está
condicionado a la capacidad del modelo para explotarla.

En términos \emph{agregados}, ese balance desigual se salda en tablas: el $\kappa$
medio de B1 ($+0,028$) queda incluso por debajo del control B0 ($+0,037$), y a un
coste de inferencia 2,2 veces mayor. Conviene ser explícito al respecto. La
evaluación mide la arquitectura, no la defiende: en media, y sobre estos tres
modelos, empaquetar la metodología como \textit{Agent Skills} no mejora el
acuerdo con las anotadoras expertas.

\begin{figure}[H]
    \centering
    \includegraphics[width=\textwidth]{fig_b0b1_metricas.pdf}
    \caption{Comparación B0 (\textit{baseline}) frente a B1 (Agent Skills) en las
    tres métricas, para los tres modelos ($N=1.313$). La exactitud mejora con B1
    en dos de tres modelos, pero el coeficiente $\kappa$ solo mejora en
    \texttt{gemini-3.1-flash-lite}.}
    \label{fig:iris-b0b1-metricas}
\end{figure}

\subsection{Divergencia entre exactitud y \texorpdfstring{$\kappa$}{kappa}}
\label{subsec:iris_divergencia}

La Tabla~\ref{tab:iris-b0b1} revela un fenómeno que constituye, a juicio de
este trabajo, su hallazgo metodológico más relevante: \textbf{la exactitud y el
coeficiente $\kappa$ apuntan en direcciones opuestas}. Considerando las medias,
el nivel B1 mejora la exactitud en $+0,036$ mientras empeora $\kappa$ en
$-0,009$. La Figura~\ref{fig:iris-divergencia} lo visualiza: al pasar de B0 a B1,
las flechas se desplazan a la derecha (más exactitud) pero descienden o se
mantienen en altura ($\kappa$ que no mejora), salvo en \texttt{gemini}.

\begin{figure}[H]
    \centering
    \includegraphics[width=0.72\textwidth]{fig_divergencia.pdf}
    \caption{Trayectoria de cada modelo de B0 a B1 en el plano
    exactitud--$\kappa$. Un desplazamiento hacia la derecha indica más exactitud;
    hacia arriba, más acuerdo. En dos de tres modelos la exactitud sube mientras
    $\kappa$ baja, señal de que la mejora aparente es un artefacto de la
    distribución marginal y no de la capacidad discriminativa.}
    \label{fig:iris-divergencia}
\end{figure}

El mecanismo subyacente es identificable. Las \textit{skills} vuelven al modelo
menos conservador: en \texttt{gpt-4o-mini} la proporción de respuestas «Sí»
para V25 pasa del 9\,\% en B0 al 46\,\% en B1. Dado que la prevalencia real de
esa variable es del 81\,\%, responder afirmativamente con mayor frecuencia
incrementa mecánicamente los aciertos por efecto de la distribución marginal,
sin que ello implique una mejora en la capacidad discriminativa. La exactitud
recompensa ese desplazamiento; $\kappa$, que descuenta el acuerdo esperable por
azar, no lo hace.

La consecuencia práctica es directa: \emph{un trabajo que reportase únicamente
exactitud concluiría que la arquitectura de Agent Skills mejora el sistema, y
tal conclusión sería incorrecta}. Este mismo patrón se observó de forma
independiente en las tres ablaciones descritas en la
Sección~\ref{subsec:iris_ablaciones}, lo que refuerza la recomendación de
reportar siempre exactitud, $\kappa$ y prevalencia de forma conjunta en tareas
de clasificación desbalanceadas.

\subsection{Diagnóstico por variable: análisis de \textit{lift}}
\label{subsec:iris_lift}

Para determinar si los sistemas poseen capacidad discriminativa real se emplea
el \textit{lift}, definido como el cociente entre la precisión sobre la clase
positiva y la prevalencia de dicha clase. Un valor de 1,0 indica que las
predicciones positivas del modelo no son mejores que las que produciría una
asignación aleatoria con la misma tasa de disparo; valores superiores indican
discriminación efectiva. La Tabla~\ref{tab:iris-lift} recoge estos valores para
el nivel B1.

\begin{table}[H]
    \centering
    \begingroup
    \small
    \ttabbox[\FBwidth]{
        \caption{Exactitud, $\kappa$ y \textit{lift} por variable en el nivel B1 ($N=1.313$)}
        \label{tab:iris-lift}
    }{
        \begin{tabularx}{\textwidth}{
            X
            >{\centering\arraybackslash}X
            >{\centering\arraybackslash}X
            >{\centering\arraybackslash}X
            >{\centering\arraybackslash}X
        }
            \hline
            \textbf{Variable} & \textbf{Prev.} & \textbf{Gemini} & \textbf{GPT-4o-mini} & \textbf{GPT-5.4-nano} \\
             & & \multicolumn{3}{c}{\small exactitud / $\kappa$ / \textit{lift}} \\
            \hline
            \texttt{lenguaje\_sexista} & 0,810 & 0,226 / $+0,015$ / 1,19 & 0,457 / $-0,037$ / 0,97 & 0,191 / $+0,000$ / 1,24 \\
            \texttt{masc\_generico} & 0,810 & \textbf{0,708 / $+0,261$ / 1,10} & 0,332 / $+0,034$ / 1,08 & 0,262 / $+0,016$ / 1,08 \\
            \texttt{sexismo\_discurso} & 0,428 & 0,569 / $-0,000$ / 0,99 & 0,561 / $+0,011$ / 1,05 & 0,558 / $+0,004$ / 1,02 \\
            \texttt{asimetria\_mujer\_hombre} & 0,062 & 0,909 / $+0,045$ / 1,91 & 0,813 / $+0,003$ / 1,03 & 0,931 / $-0,012$ / 0,00 \\
            \texttt{denominacion\_sexualizada} & 0,100 & 0,894 / $+0,011$ / 1,67 & 0,897 / $+0,061$ / 3,76 & 0,901 / $+0,014$ / 10,02 \\
            \hline
        \end{tabularx}
    }
    \endgroup
\end{table}

El diagnóstico permite distinguir tres modos de fallo cualitativamente
distintos, que explican por qué ninguna de las intervenciones ensayadas logró
elevar el acuerdo:

\begin{enumerate}
    \item \textbf{Ausencia de señal discriminativa en V25.} Con un \textit{lift}
    de 0,97 en \texttt{gpt-4o-mini}, las predicciones positivas del sistema no
    superan al azar. Resulta significativo que la prevalencia experta de esta
    variable alcance el 81\,\%, lo que sugiere que las anotadoras aplicaron un
    criterio considerablemente más amplio que la regla de inversión codificada
    en la metodología operativa. Se trata, por tanto, de un problema de validez
    de constructo antes que de capacidad del modelo: sistema y anotadoras no
    están respondiendo a la misma pregunta.

    \item \textbf{Infradetección generalizada.} Todas las variables disparan muy
    por debajo de su prevalencia real. Entre las causas plausibles se encuentra
    la exigencia, impuesta por la propia arquitectura, de citar evidencias
    literales del texto: ante la imposibilidad de reproducir un fragmento exacto,
    el agente tiende a resolver en negativo.

    \item \textbf{Desbalance extremo en V33 y V35.} Los valores elevados de
    \textit{lift} en V35 (3,76 y 10,02) descansan sobre un número muy reducido de
    predicciones positivas y deben interpretarse con cautela. En este régimen,
    $\kappa$ penaliza severamente incluso a clasificadores razonables, hasta el
    punto de que una exactitud de 0,90 puede coexistir con $\kappa = 0,000$.
\end{enumerate}

Únicamente \texttt{masc\_generico} bajo \texttt{gemini-3.1-flash-lite} alcanza
un acuerdo apreciable ($\kappa = 0,261$), siendo además el único caso en que
exactitud y $\kappa$ mejoran simultáneamente respecto al \textit{baseline}, lo
que descarta el artefacto de marginales y confirma una ganancia genuina.

\subsection{Ablaciones}
\label{subsec:iris_ablaciones}

Con objeto de descartar que el bajo rendimiento obedeciese a decisiones de
configuración, se ejecutaron tres ablaciones controladas sobre
\texttt{gpt-4o-mini}, modificando en cada caso un único factor. La
Tabla~\ref{tab:iris-ablaciones} sintetiza sus resultados.

\begin{table}[H]
    \centering
    \begingroup
    \small
    \ttabbox[\FBwidth]{
        \caption{Ablaciones controladas sobre \texttt{gpt-4o-mini}}
        \label{tab:iris-ablaciones}
    }{
        \begin{tabularx}{\textwidth}{
            X
            X
            >{\centering\arraybackslash}X
        }
            \hline
            \textbf{Factor} & \textbf{Efecto observado} & \textbf{$\Delta$ coste} \\
            \hline
            \textit{Prompt caching} &
            Sin efecto sobre $\kappa$. Altera el comportamiento del agente:
            al reestructurar el \textit{prompt} en roles \textit{system}/\textit{user},
            duplica el uso de herramientas. &
            $+87\,\%$ \\
            \hline
            \textit{Skills}-resumen de guías &
            Sin efecto claro sobre $\kappa$. &
            $-20\,\%$ \\
            \hline
            Umbral de detección (estricto \textit{vs.} permisivo) &
            Incrementa la exactitud (0,644 $\rightarrow$ 0,688) \emph{destruyendo}
            señal: el \textit{lift} de \texttt{masc\_generico} cae de 1,25 a 1,01.
            $\kappa$ no mejora de forma real. &
            --- \\
            \hline
        \end{tabularx}
    }
    \endgroup
\end{table}

La ablación del umbral merece comentario específico por su relevancia
metodológica. El régimen permisivo mejora la exactitud de forma apreciable,
resultado que en trabajos previos del proyecto se interpretó como una mejora
del sistema. El análisis de \textit{lift} revela, sin embargo, que la
precisión de \texttt{masc\_generico} desciende hasta igualar exactamente su
prevalencia: el clasificador deja de discriminar y pasa a acertar por
distribución. Se reproduce así, por tercera vez y de manera independiente, el
fenómeno descrito en la Sección~\ref{subsec:iris_divergencia}.

En cuanto al \textit{prompt caching}, su efecto resultó contraintuitivo. Aunque
la técnica reduce en un 33\,\% el coste unitario por \textit{token}, la
reestructuración del \textit{prompt} que exige duplica el volumen total
consumido, de modo que el coste neto se incrementa. Cabe señalar que la
implementación de esta técnica no es neutra respecto al comportamiento del
agente, extremo que rara vez se documenta en la literatura aplicada.

\subsection{Análisis de errores y validez de la tarea}
\label{subsec:iris_errores}

El bajo acuerdo absoluto observado en las secciones anteriores plantea una
pregunta que condiciona la interpretación de todo el capítulo: ¿refleja una
limitación de los modelos, o un límite intrínseco de la propia tarea? Esta
sección aborda esa cuestión mediante tres análisis complementarios.

\subsubsection{Heterogeneidad entre anotadores}

El corpus registra, para cada pieza, la persona que la codificó. Las 1.313
noticias anotadas se reparten entre dos equipos —Indexa (548 piezas) y UCM3
(765)—. Aunque la ausencia de doble codificación
(Sección~\ref{subsec:limitacion_fiabilidad}) impide calcular la fiabilidad
intercodificadora, sí es posible medir el acuerdo del sistema con cada equipo
por separado, a modo de análisis de sensibilidad. La
Tabla~\ref{tab:iris-sensibilidad} recoge el resultado para el nivel B1.

\begin{table}[H]
    \centering
    \begingroup
    \small
    \ttabbox[\FBwidth]{
        \caption{Análisis de sensibilidad: acuerdo (B1) según el equipo de anotación}
        \label{tab:iris-sensibilidad}
    }{
        \begin{tabularx}{\textwidth}{
            X
            X
            >{\centering\arraybackslash}X
            >{\centering\arraybackslash}X
            >{\columncolor{gray!15}\centering\arraybackslash}X
        }
            \hline
            \textbf{Modelo} & \textbf{Subconjunto} & \textbf{Exactitud} & \textbf{F1 macro} & \textbf{$\kappa$} \\
            \hline
            \multirow{3}{*}{\texttt{gemini-3.1-flash-lite}}
              & Todas (1.313)   & 0,661 & 0,443 & $+0,066$ \\
              & Solo Indexa (548) & \textbf{0,741} & \textbf{0,503} & $\mathbf{+0,118}$ \\
              & Solo UCM3 (765)  & 0,604 & 0,398 & $+0,043$ \\
            \hline
            \multirow{3}{*}{\texttt{gpt-4o-mini}}
              & Todas           & 0,612 & 0,438 & $+0,014$ \\
              & Solo Indexa     & 0,693 & 0,482 & $+0,042$ \\
              & Solo UCM3       & 0,554 & 0,398 & $+0,001$ \\
            \hline
            \multirow{3}{*}{\texttt{gpt-5.4-nano}}
              & Todas           & 0,569 & 0,364 & $+0,004$ \\
              & Solo Indexa     & 0,661 & 0,406 & $+0,014$ \\
              & Solo UCM3       & 0,502 & 0,323 & $-0,001$ \\
            \hline
        \end{tabularx}
    }
    \endgroup
\end{table}

El acuerdo con Indexa es sistemáticamente muy superior al obtenido con UCM3, en
los tres modelos y en las tres métricas de forma coherente. Para
\texttt{gemini-3.1-flash-lite}, restringir la evaluación a Indexa prácticamente
duplica el coeficiente $\kappa$ (de $0,066$ a $0,118$) y eleva la exactitud
ocho puntos; contra UCM3 en solitario, $\kappa$ desciende hacia el azar (hasta
valores negativos en \texttt{gpt-5.4-nano}). \textbf{La magnitud de esta
variabilidad entre equipos es comparable a la diferencia entre modelos}, lo que
indica que el origen de la anotación pesa en el acuerdo tanto como la elección
del modelo.

Debe subrayarse que este resultado \emph{no} autoriza a reportar únicamente el
subconjunto más favorable: al no existir doble codificación, no puede
determinarse qué equipo refleja mejor el constructo, y la diferencia está además
confundida con el hecho de que cada equipo codificó piezas distintas. El
análisis acota el techo alcanzable, pero no lo atribuye a un error de anotación.

\subsubsection{Casos de consenso de los modelos frente al \textit{ground truth}}

Un segundo análisis examina las celdas en las que \textbf{los tres modelos
coinciden entre sí} pero difieren de la anotación experta: son las más
diagnósticas, pues descartan el error idiosincrásico de un modelo concreto. De
las 6.565 celdas evaluadas (1.313 piezas $\times$ 5 variables), 1.439 presentan
consenso de los tres modelos en contra del \textit{ground truth}. Su dirección
es prácticamente unánime: en \textbf{1.436 casos} los tres modelos responden «No»
donde la anotación dice «Sí», y solo en 3 casos ocurre lo contrario. El
desacuerdo es, por tanto, un \textbf{sesgo sistemático de infradetección}: los
modelos detectan de forma consistente menos sexismo que las anotadoras, y no de
manera aleatoria.

Se descartó empíricamente que este fenómeno obedeciese a defectos en la
extracción del texto: solo 49 de las 1.313 piezas (3,7\,\%) presentan
contenido sospechoso —muy corto o compuesto por texto auxiliar de la web
(avisos de suscripción, fichas de autoría)— y su exclusión no altera las
métricas (para \texttt{gemini-3.1-flash-lite}, la exactitud pasa de $0,661$ a
$0,660$). El desacuerdo se produce, por tanto, sobre artículos completos y
correctamente extraídos.

\subsubsection{Lectura cualitativa}

La revisión manual de casos de consenso permite atribuir el desacuerdo a tres
mecanismos, ninguno de los cuales constituye un fallo de detección en sentido
estricto:

\begin{itemize}
    \item \textbf{Contradicción interna del codebook (V25).} La definición de
    \texttt{lenguaje\_sexista} combina una regla estricta —la regla de inversión—
    con una cláusula agregadora según la cual todo subtipo de sexismo lingüístico
    implica marcar V25. Los modelos aplican la primera; la anotación, a menudo,
    la segunda. Ambas lecturas son defendibles porque el propio criterio es
    ambiguo.

    \item \textbf{Umbral del masculino genérico (V26).} El español emplea el
    masculino genérico de forma pervasiva (\textit{los trabajadores}, \textit{los
    ciudadanos}). La anotación tiende a marcarlo sistemáticamente —lo que explica
    su prevalencia del 81\,\%—, mientras que los modelos aplican un umbral más
    alto y solo lo señalan cuando resulta evitable e invisibilizador. Sistema y
    anotadoras responden, de hecho, a preguntas distintas.

    \item \textbf{Anotación sobre-inclusiva (V30).} En una fracción de los casos
    de \texttt{sexismo\_discurso} no se aprecia, en una lectura atenta del texto,
    el fenómeno codificado; estos casos se concentran en el anotador de menor
    acuerdo global.
\end{itemize}

En conjunto, estos análisis reencuadran el resultado principal: la baja
concordancia \textbf{no debe leerse como incapacidad de los modelos para detectar
sexismo}, sino como el producto de un criterio experto notablemente amplio, de
una heterogeneidad sustancial entre anotadores y de la ambigüedad de algunas
definiciones del codebook. El coeficiente $\kappa$ mide, en buena medida, esa
distancia de criterio.

\subsection{Coste computacional}
\label{subsec:iris_coste}

La Tabla~\ref{tab:iris-b0b1} incluye el coste real de inferencia. La
arquitectura de agentes multiplica por 2,2 el coste respecto al \textit{baseline}
(30,91 USD frente a 13,83 USD para los tres modelos sobre 1.313 noticias), al
requerir aproximadamente diez llamadas por variable frente a una única llamada.

Esta diferencia condiciona la viabilidad de la aproximación con modelos
desplegados localmente. Las mediciones efectuadas sobre \texttt{gemma4:e4b}
arrojan un tiempo aproximado de ocho minutos por artículo, lo que situaría el
procesamiento del corpus de evaluación en torno a las 175 horas. El coste
computacional de las arquitecturas agénticas constituye, por tanto, una
limitación estructural para su despliegue fuera de infraestructura dedicada.
